\documentclass[../main.tex]{subfiles}

\title{1.1 - Limits}
\showsolutionsfalse

\begin{document}

\maketitle

\goals{Describe limits graphically, numerically, and algebraically.}

\beforeclassvideos{\href{https://web.uvic.ca/~tbazett/Calculus/sec-intro-limits.html}{1.1.1 and 1.1.2}}

\vspace{0.25em}
\hrule
\vspace{0.5em}
\
\Example Consider these three functions:

\[
f(x)=x+1,
\qquad
g(x)=\frac{x^2-1}{x-1}=\frac{(x+1)(x-1)}{x-1}, \qquad
h(x)=
\begin{cases}
x+1,&x\ne1,\\
3,&x=1.
\end{cases}
\]

 How are these similar or different?

\sol{
\begin{center}
\begin{tikzpicture}[x=0.9cm,y=0.7cm]
  \foreach \s in {0,4,8}{
    \begin{scope}[xshift=\s cm]
      \draw[note axis] (-0.5,0)--(2.8,0);
      \draw[note axis] (0,-0.2)--(0,2.8);
      \draw[thick] (-0.3,0.7)--(2.2,3.2);
    \end{scope}
  }
  \node[above] at (1.7,2.3){$f(x)$};
  \node[open point] at (5,2){};
  \draw[note dashed] (5,0)--(5,2);
  \node[above] at (5.8,2.3){$g(x)$};
  \node[open point] at (9,2){};
  \node[closed point] at (9,2.65){};
  \node[above] at (9.8,2.3){$h(x)$};
\end{tikzpicture}
\end{center}

\[
\text{Domain:}\qquad
\mathbb{R},\qquad(-\infty,1)\cup(1,\infty),\qquad\mathbb{R}.
\]

Same everywhere but $x=1$, where $f(1)=2$, $g(1)$ undefined, $h(1)=3$.
}
\Definition{$\displaystyle\lim_{x\to a}f(x)=L$ if $f(x)$ can be made arbitrarily close to $L$ for all values of $x$ sufficiently close to $a$.}

\Example Compute the limit as $x\to 2$ of $f(x)$, $g(x)$, and $h(x)$. 

\sol{
\[
\lim_{x\to1}f(x)=\lim_{x\to1}g(x)=\lim_{x\to1}h(x)=2
\]

\vspace{1in}}

\Example Use a table of values to estimate $\displaystyle\lim_{x\to1}(x+1)$. 

\sol{
\begin{center}
\begin{tabular}{c|c}
$x$&$f(x)$\\ \hline
0.9&1.9\\
0.99&1.99\\
0.999&1.999\\
1.1&2.1\\
1.01&2.01\\
1.001&2.001
\end{tabular}
\qquad slope? closer to 2.
\end{center}
}

\newpage

\Example Use a table of values to estimate $\ds\lim_{x\to0}\frac1x$. 

\sol{

\begin{center}
\begin{tabular}{c|c}
$x$&$f(x)$\\ \hline
0.1&10\\
0.01&100\\
0.001&1000
\end{tabular}
\qquad not getting close to anything!
\end{center}

So limit does not exist!

\begin{center}
\begin{tikzpicture}[x=0.75cm,y=0.45cm]
  \draw[note axis] (-4,0)--(4,0);
  \draw[note axis] (0,-5)--(0,5);
  \draw[note curve,domain=0.22:4,samples=80] plot (\x,{1/\x});
  \draw[note curve,domain=-4:-0.22,samples=80] plot (\x,{1/\x});
\end{tikzpicture}
\end{center}
}

\Example Investigate
\[
f(x)=
\begin{cases}
x,&x<1,\\
x+1,&x\ge1.
\end{cases}
\]

\sol{

\begin{center}
\begin{tikzpicture}[x=1cm,y=0.8cm]
  \draw[note axis] (-1,0)--(3,0);
  \draw[note axis] (0,-0.5)--(0,3.5);
  \draw[thick] (-0.7,-0.7)--(1,1);
  \draw[thick] (1,2)--(2.3,3.3);
  \node[open point] at (1,1){};
  \node[closed point] at (1,2){};
\end{tikzpicture}
\end{center}}

\Definition{$\displaystyle\lim_{x\to a^+}f(x)$ means $f(x)$ gets arbitrarily close to $L$ for all $x$ sufficiently close to $a$ from the right, i.e. $x>a$.}

\sol{
\[
\lim_{x\to1^+}f(x)=
\qquad
\lim_{x\to1^-}f(x)=
\]}
\newpage 
\Example Investigate limits of 


\[
f(x)=\sin\!\left(\frac1x\right)\]
  
\sol{
\begin{center}
\begin{tikzpicture}[x=0.52cm,y=1.1cm]
  \draw[note axis] (0,0)--(14,0);
  \draw[note axis] (0,-1.6)--(0,1.6);
  \draw[note curve,domain=0.13:14,samples=300] plot (\x,{sin((1/\x) r)});
\end{tikzpicture}
\end{center}
\[
\lim_{x\to0^+}f(x)\ \DNE,\quad\text{as }f(x)\text{ gets close to every }y\text{ value in }[-1,1].
\]}

\vfill
\vspace{0.25em}
\hrule
\vspace{0.5em}

\Example Graphically investigate limits of $\frac{x^2-1}{x^2-x-2}$

\sol{\vspace{3in}}

\afterclassproblems{Section 1.1 problems 9,11,15}


\end{document}
